\documentclass[11pt]{article}
\usepackage[margin=1in]{geometry}
\usepackage{amsmath, amssymb, amsthm}
\usepackage{hyperref}
\usepackage{enumitem}

\title{Homework Practice: Exact Differential Equations}
\author{Math Sections Curator}
\date{September 2026}

\begin{document}

\maketitle
\tableofcontents
\newpage

% ---------------------------------------------------------
% SECTION 1: Concept Checklist & Question Mapping
% ---------------------------------------------------------
\section{Concept Checklist \& Question Mapping}

To fully master the concepts of Exact Differential Equations as outlined in the reference material, you must be comfortable with the following problem types, techniques, and common pitfalls. This checklist links each concept to its corresponding practice problems.

\begin{itemize}
    \item \textbf{Concept 1: Classifying Differential Equations} (Questions 1--4)
    \begin{itemize}
        \item Testing for Separable, Linear, and Exact forms.
    \end{itemize}
    \item \textbf{Concept 2: The Exactness Test \& Minus Sign Traps} (Questions 5--8)
    \begin{itemize}
        \item Verifying $M_y = N_x$.
        \item Distributing negative signs before defining $N(x,y)$. 
    \end{itemize}
    \item \textbf{Concept 3: Standard Exact Equations} (Questions 9--14)
    \begin{itemize}
        \item Integrating $M$ with respect to $x$ and solving for $g(y)$.
        \item \textit{Calculus Trick:} Integration by Parts (Question 11).
        \item \textit{Calculus Trick:} Fractional Exponents (Question 12).
        \item \textit{Calculus Trick:} Logarithms and Absolute Values (Question 13).
    \end{itemize}
    \item \textbf{Concept 4: Initial Value Problems (IVPs)} (Questions 15--20)
    \begin{itemize}
        \item Solving for the specific constant $C$.
        \item \textit{Trap:} Dropping absolute value bars appropriately (Question 18).
        \item \textit{Trap:} Choosing the correct explicit sign branch (Question 17).
    \end{itemize}
    \item \textbf{Concept 5: Integrating Factors (Non-Exact to Exact)} (Questions 21--25)
    \begin{itemize}
        \item Using $\mu(x) = e^{\int \frac{M_y - N_x}{N} dx}$.
        \item Using $\mu(y) = e^{\int \frac{N_x - M_y}{M} dy}$.
    \end{itemize}
    \item \textbf{Concept 6: Diagnostic Testing ("Find the Flaw")} (Questions 26--30)
    \begin{itemize}
        \item Spotting double-counting of integration terms.
        \item Spotting variable mix-ups ($M_x$ vs $M_y$).
        \item Spotting integration amnesia ($g(y)$ vs $g'(y)$).
        \item Spotting the expression trap (Forgetting $= C$).
    \end{itemize}
\end{itemize}

\newpage

% ---------------------------------------------------------
% SECTION 2: Problem Sets
% ---------------------------------------------------------
\section{Problem Sets}

\subsection*{Part 1: Classification}
Determine if the following equations are Separable, Linear, Exact, or none of these. Note that equations can have more than one classification.
\begin{enumerate}[label=\textbf{\arabic*.}]
    \item $(y - x^2)dx + x dy = 0$
    \item $x^2 dx + y^2 dy = 0$
    \item $\sqrt{x-y} dx + dy = 0$
    \item $x \frac{dy}{dx} - 3y = x^4$
\end{enumerate}

\subsection*{Part 2: The Exactness Test}
Determine whether the following equations are exact. Do not solve them.
\begin{enumerate}[label=\textbf{\arabic*.}, resume]
    \item $(4x^3 y - 15x^2 - y)dx + (x^4 + 3y^2 - x)dy = 0$
    \item $(\sin(y) - y \sin(x))dx + (x \cos(y) + \cos(x))dy = 0$
    \item $2xydx - (x^2 + y^2)dy = 0$
    \item $(e^x \cos y)dx + (e^x \sin y)dy = 0$
\end{enumerate}

\subsection*{Part 3: Standard Exact Equations}
Determine whether the equations are exact. If exact, find the general implicit solution.
\begin{enumerate}[label=\textbf{\arabic*.}, resume]
    \item $(3x^2 y^2 - x)dx + (2x^3 y + y)dy = 0$
    \item $(y \cos(x) + 2xe^y)dx + (\sin(x) + x^2 e^y - 1)dy = 0$
    \item $(y e^x - x e^x)dx + (e^x)dy = 0$
    \item $(\frac{2}{3}x^{-1/3} y + 1)dx + (x^{2/3} + y^{-1/2})dy = 0$
    \item $(\frac{1}{x} + y^2)dx + (2xy + \frac{1}{y})dy = 0$
    \item $(\frac{y}{x^2+1})dx + (\arctan(x) + 2y)dy = 0$
\end{enumerate}

\subsection*{Part 4: Initial Value Problems (IVPs)}
Solve the following initial value problems.
\begin{enumerate}[label=\textbf{\arabic*.}, resume]
    \item $(2x - y)dx + (2y - x)dy = 0, \quad y(1) = 3$
    \item $(\cos(x)\cos(y))dx - (\sin(x)\sin(y))dy = 0, \quad y(\pi/2) = 0$
    \item Solve explicitly for $y$: $2x dx + 2y dy = 0, \quad y(0) = -5$
    \item $(\frac{1}{x})dx + (\frac{1}{y})dy = 0, \quad y(1) = 2$
    \item $(e^y)dx + (x e^y - \sin(y))dy = 0, \quad y(2) = 0$
    \item $(\frac{2x}{y})dx - (\frac{x^2}{y^2})dy = 0, \quad y(1) = 1$
\end{enumerate}

\subsection*{Part 5: Integrating Factors}
The following equations are NOT exact. Find the appropriate integrating factor ($\mu(x)$ or $\mu(y)$), multiply the equation by it, and solve the newly exact equation.
\begin{enumerate}[label=\textbf{\arabic*.}, resume]
    \item $(2y^2 + 3x)dx + (2xy)dy = 0$
    \item $(2xy^3)dx + (3x^2 y^2 + x^2 y^3)dy = 0$
    \item $(4xy + 3y^2 - x)dx + (x^2 + 2xy)dy = 0$
    \item $(y)dx + (2x - y e^y)dy = 0$
    \item $(y)dx + (-x + y^2)dy = 0$
\end{enumerate}

\subsection*{Part 6: Diagnostic Testing ("Find the Flaw")}
Identify the mathematical flaw in the provided student solutions, explain the error, and state the correct step/solution.
\begin{enumerate}[label=\textbf{\arabic*.}, resume]
    \item \textbf{Equation:} $(3x^2 y^2)dx + (2x^3 y)dy = 0$ \\
    \textbf{Flawed Solution:} $M_y = 6x^2 y, N_x = 6x^2 y$. Exact. $\int M dx = x^3 y^2$. $\int N dy = x^3 y^2$. Combine terms: $F(x,y) = 2x^3 y^2 = C$.
    
    \item \textbf{Equation:} $2xydx - (x^2)dy = 0$ \\
    \textbf{Flawed Solution:} $M = 2xy \implies M_y = 2x$. $N = x^2 \implies N_x = 2x$. Since $2x = 2x$, the equation is exact.
    
    \item \textbf{Equation:} $(2x+y)dx + (x+2y)dy = 0$ \\
    \textbf{Flawed Solution:} Let $M = 2x+y$ and $N = x+2y$. Taking derivatives, $M_x = 2$ and $N_y = 2$. Therefore, the equation is exact.
    
    \item \textbf{Equation:} $(y^2)dx + (2xy + \sin y)dy = 0$ \\
    \textbf{Flawed Solution:} $M_y = 2y, N_x = 2y$. Exact. $F(x,y) = \int y^2 dx = xy^2 + g(y)$. Differentiate with respect to $y$: $F_y = 2xy + g(y)$. Set equal to $N$: $2xy + g(y) = 2xy + \sin y \implies g(y) = \sin y$. Solution is $xy^2 + \sin y = C$.
    
    \item \textbf{Equation:} $(1)dx + (1)dy = 0$ \\
    \textbf{Flawed Solution:} $M_y = 0, N_x = 0$. Exact. $F(x,y) = \int 1 dx = x + g(y)$. $F_y = g'(y) = 1 \implies g(y) = y$. The solution is $x + y$.
\end{enumerate}

\newpage

% ---------------------------------------------------------
% SECTION 3: Detailed Solutions
% ---------------------------------------------------------
\section{Detailed Solutions}

\subsection*{Part 1: Classification}

\textbf{1. $(y - x^2)dx + x dy = 0$} \\
\textbf{Exactness:} $M = y - x^2 \implies M_y = 1$. $N = x \implies N_x = 1$. It is \textbf{Exact}. \\
\textbf{Linearity:} Divide by $xdx$: $\frac{y}{x} - x + \frac{dy}{dx} = 0 \implies \frac{dy}{dx} + \frac{1}{x}y = x$. It matches the form $y' + P(x)y = Q(x)$. It is \textbf{Linear}. \\
\textbf{Separability:} Cannot factor $y - x^2$ into $f(x)g(y)$. Not separable. \\
\textbf{Final Answer:} Linear and Exact.

\textbf{2. $x^2 dx + y^2 dy = 0$} \\
\textbf{Exactness:} $M_y = 0, N_x = 0$. It is \textbf{Exact}. \\
\textbf{Linearity:} $y^2$ means it is not linear in $y$. \\
\textbf{Separability:} Can be written as $y^2 dy = -x^2 dx$. It is \textbf{Separable}. \\
\textbf{Final Answer:} Exact and Separable.

\textbf{3. $\sqrt{x-y} dx + dy = 0$} \\
\textbf{Exactness:} $M_y = -\frac{1}{2\sqrt{x-y}} \neq 0 = N_x$. Not exact. \\
\textbf{Linearity:} $y$ is trapped in a square root. Not linear. \\
\textbf{Separability:} Cannot separate $x$ and $y$ from the square root. Not separable. \\
\textbf{Final Answer:} None of these.

\textbf{4. $x \frac{dy}{dx} - 3y = x^4$} \\
\textbf{Linearity:} Divide by $x$: $\frac{dy}{dx} - \frac{3}{x}y = x^3$. It is \textbf{Linear}. \\
\textbf{Exactness:} Rearranging gives $-(x^4 + 3y)dx + x dy = 0$. $M_y = -3, N_x = 1$. Not exact. \\
\textbf{Separability:} $\frac{dy}{dx} = x^3 + \frac{3y}{x}$. Cannot separate $x$ and $y$ via multiplication. Not separable. \\
\textbf{Final Answer:} Linear only.

\subsection*{Part 2: The Exactness Test}

\textbf{5. $(4x^3 y - 15x^2 - y)dx + (x^4 + 3y^2 - x)dy = 0$} \\
$M_y = \frac{\partial}{\partial y}(4x^3 y - 15x^2 - y) = 4x^3 - 1$ \\
$N_x = \frac{\partial}{\partial x}(x^4 + 3y^2 - x) = 4x^3 - 1$ \\
$M_y = N_x$. \textbf{Exact.}

\textbf{6. $(\sin(y) - y \sin(x))dx + (x \cos(y) + \cos(x))dy = 0$} \\
$M_y = \cos(y) - \sin(x)$ \\
$N_x = \cos(y) - \sin(x)$ \\
$M_y = N_x$. \textbf{Exact.}

\textbf{7. $2xydx - (x^2 + y^2)dy = 0$} \\
Remember to distribute the minus sign! $N(x,y) = -x^2 - y^2$. \\
$M_y = 2x$ \\
$N_x = -2x$ \\
$M_y \neq N_x$. \textbf{Not exact.}

\textbf{8. $(e^x \cos y)dx + (e^x \sin y)dy = 0$} \\
$M_y = -e^x \sin y$ \\
$N_x = e^x \sin y$ \\
$M_y \neq N_x$. \textbf{Not exact.}

\subsection*{Part 3: Standard Exact Equations}

\textbf{9. $(3x^2 y^2 - x)dx + (2x^3 y + y)dy = 0$} \\
\textbf{Step 1:} $M_y = 6x^2 y, N_x = 6x^2 y$. Exact. \\
\textbf{Step 2:} $F(x,y) = \int M dx = \int (3x^2 y^2 - x) dx = x^3 y^2 - \frac{1}{2}x^2 + g(y)$ \\
\textbf{Step 3:} $F_y = 2x^3 y + g'(y)$. Set to $N$: $2x^3 y + g'(y) = 2x^3 y + y \implies g'(y) = y$. \\
\textbf{Step 4:} $g(y) = \frac{1}{2}y^2$. \\
\textbf{Final Answer:} $x^3 y^2 - \frac{1}{2}x^2 + \frac{1}{2}y^2 = C$

\textbf{10. $(y \cos(x) + 2xe^y)dx + (\sin(x) + x^2 e^y - 1)dy = 0$} \\
\textbf{Step 1:} $M_y = \cos(x) + 2xe^y$. $N_x = \cos(x) + 2xe^y$. Exact. \\
\textbf{Step 2:} $F(x,y) = \int M dx = y \sin(x) + x^2 e^y + g(y)$ \\
\textbf{Step 3:} $F_y = \sin(x) + x^2 e^y + g'(y) = \sin(x) + x^2 e^y - 1 \implies g'(y) = -1$. \\
\textbf{Step 4:} $g(y) = -y$. \\
\textbf{Final Answer:} $y \sin(x) + x^2 e^y - y = C$

\textbf{11. $(y e^x - x e^x)dx + (e^x)dy = 0$} \\
\textbf{Step 1:} $M_y = e^x$, $N_x = e^x$. Exact. \\
\textbf{Step 2:} Integrating $N$ with respect to $y$ is easier! $F(x,y) = \int e^x dy = y e^x + h(x)$. \\
\textbf{Step 3:} $F_x = y e^x + h'(x) = y e^x - x e^x \implies h'(x) = -x e^x$. \\
\textbf{Step 4:} Integration by parts for $\int -x e^x dx$. Let $u = -x, dv = e^x dx \implies du = -dx, v = e^x$. \\
$\int u dv = -xe^x - \int -e^x dx = -x e^x + e^x$. \\
\textbf{Final Answer:} $y e^x - x e^x + e^x = C$

\textbf{12. $(\frac{2}{3}x^{-1/3} y + 1)dx + (x^{2/3} + y^{-1/2})dy = 0$} \\
\textbf{Step 1:} $M_y = \frac{2}{3}x^{-1/3}$, $N_x = \frac{2}{3}x^{-1/3}$. Exact. \\
\textbf{Step 2:} $F(x,y) = \int M dx = x^{2/3} y + x + g(y)$. \\
\textbf{Step 3:} $F_y = x^{2/3} + g'(y) = x^{2/3} + y^{-1/2} \implies g'(y) = y^{-1/2}$. \\
\textbf{Step 4:} $g(y) = \int y^{-1/2} dy = 2y^{1/2} = 2\sqrt{y}$. \\
\textbf{Final Answer:} $x^{2/3} y + x + 2\sqrt{y} = C$

\textbf{13. $(\frac{1}{x} + y^2)dx + (2xy + \frac{1}{y})dy = 0$} \\
\textbf{Step 1:} $M_y = 2y$, $N_x = 2y$. Exact. \\
\textbf{Step 2:} $F(x,y) = \int (\frac{1}{x} + y^2)dx = \ln|x| + xy^2 + g(y)$. \\
\textbf{Step 3:} $F_y = 2xy + g'(y) = 2xy + \frac{1}{y} \implies g'(y) = \frac{1}{y}$. \\
\textbf{Step 4:} $g(y) = \ln|y|$. \\
\textbf{Final Answer:} $\ln|x| + xy^2 + \ln|y| = C$

\textbf{14. $(\frac{y}{x^2+1})dx + (\arctan(x) + 2y)dy = 0$} \\
\textbf{Step 1:} $M_y = \frac{1}{x^2+1}$, $N_x = \frac{1}{x^2+1}$. Exact. \\
\textbf{Step 2:} $F(x,y) = \int \frac{y}{x^2+1} dx = y \arctan(x) + g(y)$. \\
\textbf{Step 3:} $F_y = \arctan(x) + g'(y) = \arctan(x) + 2y \implies g'(y) = 2y$. \\
\textbf{Step 4:} $g(y) = y^2$. \\
\textbf{Final Answer:} $y \arctan(x) + y^2 = C$

\subsection*{Part 4: Initial Value Problems}

\textbf{15. $(2x - y)dx + (2y - x)dy = 0, \quad y(1) = 3$} \\
$M_y = -1, N_x = -1$. Exact. \\
$F(x,y) = \int (2x-y)dx = x^2 - xy + g(y)$. \\
$F_y = -x + g'(y) = 2y - x \implies g'(y) = 2y \implies g(y) = y^2$. \\
General solution: $x^2 - xy + y^2 = C$. \\
Apply IVP: $(1)^2 - (1)(3) + (3)^2 = 1 - 3 + 9 = 7$. \\
\textbf{Final Answer:} $x^2 - xy + y^2 = 7$

\textbf{16. $(\cos(x)\cos(y))dx - (\sin(x)\sin(y))dy = 0, \quad y(\pi/2) = 0$} \\
$M_y = -\cos(x)\sin(y), N_x = -\cos(x)\sin(y)$. Exact. \\
$F(x,y) = \int \cos(x)\cos(y)dx = \sin(x)\cos(y) + g(y)$. \\
$F_y = -\sin(x)\sin(y) + g'(y) = -\sin(x)\sin(y) \implies g'(y) = 0 \implies g(y) = 0$. \\
General solution: $\sin(x)\cos(y) = C$. \\
Apply IVP: $\sin(\pi/2)\cos(0) = (1)(1) = 1 \implies C = 1$. \\
\textbf{Final Answer:} $\sin(x)\cos(y) = 1$

\textbf{17. $2x dx + 2y dy = 0, \quad y(0) = -5$} \\
$M_y = 0, N_x = 0$. Exact. \\
$F(x,y) = x^2 + y^2 = C$. \\
Apply IVP: $0^2 + (-5)^2 = C \implies C = 25$. \\
Implicit: $x^2 + y^2 = 25$. Solving for $y$: $y = \pm\sqrt{25-x^2}$. \\
Since $y(0) = -5$, we must take the negative branch. \\
\textbf{Final Answer:} $y = -\sqrt{25 - x^2}$

\textbf{18. $(\frac{1}{x})dx + (\frac{1}{y})dy = 0, \quad y(1) = 2$} \\
$M_y = 0, N_x = 0$. Exact. \\
$F(x,y) = \ln|x| + \ln|y| = C$. \\
Because the initial condition $x=1, y=2$ defines a domain strictly in the first quadrant, we drop the absolute value bars: $\ln(x) + \ln(y) = C \implies \ln(xy) = C$. \\
Apply IVP: $\ln(1 \cdot 2) = C \implies C = \ln(2)$. \\
$\ln(xy) = \ln(2) \implies xy = 2$. \\
\textbf{Final Answer:} $xy = 2$ or $y = \frac{2}{x}$

\textbf{19. $(e^y)dx + (x e^y - \sin(y))dy = 0, \quad y(2) = 0$} \\
$M_y = e^y, N_x = e^y$. Exact. \\
$F(x,y) = \int e^y dx = x e^y + g(y)$. \\
$F_y = x e^y + g'(y) = x e^y - \sin(y) \implies g'(y) = -\sin(y) \implies g(y) = \cos(y)$. \\
General Solution: $x e^y + \cos(y) = C$. \\
Apply IVP: $2 e^0 + \cos(0) = C \implies 2(1) + 1 = 3 \implies C = 3$. \\
\textbf{Final Answer:} $x e^y + \cos(y) = 3$

\textbf{20. $(\frac{2x}{y})dx - (\frac{x^2}{y^2})dy = 0, \quad y(1) = 1$} \\
$M_y = -\frac{2x}{y^2}, N_x = -\frac{2x}{y^2}$. Exact. \\
$F(x,y) = \int \frac{2x}{y} dx = \frac{x^2}{y} + g(y)$. \\
$F_y = -\frac{x^2}{y^2} + g'(y) = -\frac{x^2}{y^2} \implies g(y) = 0$. \\
General Solution: $\frac{x^2}{y} = C$. \\
Apply IVP: $\frac{1^2}{1} = C \implies C = 1$. \\
\textbf{Final Answer:} $\frac{x^2}{y} = 1$ (or $y = x^2$)

\subsection*{Part 5: Integrating Factors}

\textbf{21. $(2y^2 + 3x)dx + (2xy)dy = 0$} \\
$M_y = 4y, N_x = 2y$. $\frac{M_y - N_x}{N} = \frac{4y - 2y}{2xy} = \frac{2y}{2xy} = \frac{1}{x}$. \\
$\mu(x) = e^{\int \frac{1}{x} dx} = x$. \\
Multiply by $x$: $(2xy^2 + 3x^2)dx + (2x^2 y)dy = 0$. \\
Now $M_y = 4xy, N_x = 4xy$. Exact! \\
$F(x,y) = \int (2x^2 y) dy = x^2 y^2 + h(x)$. \\
$F_x = 2xy^2 + h'(x) = 2xy^2 + 3x^2 \implies h'(x) = 3x^2 \implies h(x) = x^3$. \\
\textbf{Final Answer:} $x^2 y^2 + x^3 = C$

\textbf{22. $(2xy^3)dx + (3x^2 y^2 + x^2 y^3)dy = 0$} \\
$M_y = 6xy^2, N_x = 6xy^2 + 2x y^3$. \\
$\frac{N_x - M_y}{M} = \frac{2x y^3}{2x y^3} = 1$. \\
$\mu(y) = e^{\int 1 dy} = e^y$. \\
Multiply by $e^y$: $(2x y^3 e^y)dx + x^2(3y^2 + y^3)e^y dy = 0$. \\
$F(x,y) = \int 2xy^3 e^y dx = x^2 y^3 e^y + g(y)$. \\
$F_y = x^2(3y^2 e^y + y^3 e^y) + g'(y) = N \implies g'(y) = 0$. \\
\textbf{Final Answer:} $x^2 y^3 e^y = C$

\textbf{23. $(4xy + 3y^2 - x)dx + (x^2 + 2xy)dy = 0$} \\
$M_y = 4x + 6y, N_x = 2x + 2y$. \\
$\frac{M_y - N_x}{N} = \frac{2x + 4y}{x^2 + 2xy} = \frac{2(x + 2y)}{x(x + 2y)} = \frac{2}{x}$. \\
$\mu(x) = e^{\int 2/x dx} = x^2$. \\
Multiply by $x^2$: $(4x^3 y + 3x^2 y^2 - x^3)dx + (x^4 + 2x^3 y)dy = 0$. \\
$F(x,y) = \int (x^4 + 2x^3 y)dy = x^4 y + x^3 y^2 + h(x)$. \\
$F_x = 4x^3 y + 3x^2 y^2 + h'(x) = M \implies h'(x) = -x^3 \implies h(x) = -\frac{1}{4}x^4$. \\
\textbf{Final Answer:} $x^4 y + x^3 y^2 - \frac{1}{4}x^4 = C$

\textbf{24. $(y)dx + (2x - y e^y)dy = 0$} \\
$M_y = 1, N_x = 2$. \\
$\frac{N_x - M_y}{M} = \frac{2 - 1}{y} = \frac{1}{y}$. \\
$\mu(y) = e^{\int 1/y dy} = y$. \\
Multiply by $y$: $y^2 dx + (2xy - y^2 e^y)dy = 0$. \\
$F(x,y) = \int y^2 dx = x y^2 + g(y)$. \\
$F_y = 2xy + g'(y) = 2xy - y^2 e^y \implies g'(y) = -y^2 e^y$. \\
Integration by parts (tabular method) yields $g(y) = -y^2 e^y + 2y e^y - 2e^y$. \\
\textbf{Final Answer:} $xy^2 - y^2 e^y + 2y e^y - 2e^y = C$

\textbf{25. $(y)dx + (-x + y^2)dy = 0$} \\
$M_y = 1, N_x = -1$. \\
$\frac{N_x - M_y}{M} = \frac{-1 - 1}{y} = -\frac{2}{y}$. \\
$\mu(y) = e^{\int -2/y dy} = y^{-2}$. \\
Multiply by $y^{-2}$: $\frac{1}{y} dx + (-\frac{x}{y^2} + 1)dy = 0$. \\
$F(x,y) = \int \frac{1}{y} dx = \frac{x}{y} + g(y)$. \\
$F_y = -\frac{x}{y^2} + g'(y) = -\frac{x}{y^2} + 1 \implies g'(y) = 1 \implies g(y) = y$. \\
\textbf{Final Answer:} $\frac{x}{y} + y = C$

\subsection*{Part 6: Diagnostic Testing ("Find the Flaw")}

\textbf{26.} \\
\textbf{The Flaw:} The student double-counted the shared term $x^3 y^2$ when combining the integrals of $M$ and $N$. Overlapping terms that appear in both partial integrations should only be written once in the potential function. \\
\textbf{Correct Step:} $F(x,y) = x^3 y^2 = C$.

\textbf{27.} \\
\textbf{The Flaw:} The student failed to distribute the negative sign into the $N$ term. $N$ should be $-x^2$, which makes $N_x = -2x$. Because $2x \neq -2x$, the equation is NOT exact. \\
\textbf{Correct Step:} $M_y = 2x, N_x = -2x$. Not exact.

\textbf{28.} \\
\textbf{The Flaw:} The student mixed up the variables during partial differentiation. To test for exactness, you must differentiate $M$ (the term attached to $dx$) with respect to $y$, and $N$ (attached to $dy$) with respect to $x$. \\
\textbf{Correct Step:} $M_y = 1$ and $N_x = 1$. The equation is indeed exact, but the student's justification was mathematically incorrect.

\textbf{29.} \\
\textbf{The Flaw:} The student suffered from "integration amnesia". When comparing $F_y$ to $N$, the remaining term represents $g'(y)$, not $g(y)$. The student skipped integrating $\sin y$. \\
\textbf{Correct Step:} $g'(y) = \sin y \implies g(y) = -\cos y$. The final solution is $xy^2 - \cos y = C$.

\textbf{30.} \\
\textbf{The Flaw:} The student left the final answer as an expression $x+y$. Exact differential equations yield implicit relations equal to a constant, not standalone expressions. \\
\textbf{Correct Step:} Set the function equal to $C$. The correct answer is $x + y = C$.

% ---------------------------------------------------------
% SECTION 4: Re-Printed Checklist (For Quick Reference)
% ---------------------------------------------------------
\newpage
\section{Concept Checklist \& Question Mapping (Reference Guide)}

\begin{itemize}
    \item \textbf{Concept 1: Classifying Differential Equations} $\rightarrow$ Q1, Q2, Q3, Q4
    \item \textbf{Concept 2: The Exactness Test \& Minus Sign Traps} $\rightarrow$ Q5, Q6, Q7, Q8
    \item \textbf{Concept 3: Standard Exact Equations} $\rightarrow$ Q9, Q10, Q11, Q12, Q13, Q14
    \item \textbf{Concept 4: Initial Value Problems (IVPs)} $\rightarrow$ Q15, Q16, Q17, Q18, Q19, Q20
    \item \textbf{Concept 5: Integrating Factors (Non-Exact to Exact)} $\rightarrow$ Q21, Q22, Q23, Q24, Q25
    \item \textbf{Concept 6: Diagnostic Testing ("Find the Flaw")} $\rightarrow$ Q26, Q27, Q28, Q29, Q30
\end{itemize}

\end{document}